\documentclass[11pt]{article}

\usepackage[margin=2.5cm]{geometry}
\usepackage[utf8]{inputenc}
\usepackage[T1]{fontenc}
\usepackage{amsmath,amssymb}
\usepackage{booktabs}
\usepackage{array}
\usepackage{xcolor}
\usepackage{colortbl}
\usepackage{titlesec}
\usepackage{enumitem}
\usepackage{parskip}
\usepackage{longtable}
\usepackage{fancyhdr}
\usepackage{tcolorbox}
\usepackage{graphicx}
\usepackage[colorlinks=true,linkcolor=blue!45!black,urlcolor=blue!45!black]{hyperref}

\definecolor{crit}{RGB}{155,25,25}
\definecolor{warn}{RGB}{180,110,0}
\definecolor{good}{RGB}{20,110,60}
\definecolor{hdr}{RGB}{30,55,105}

\titleformat{\section}{\normalfont\Large\bfseries\color{hdr}}{\thesection}{0.6em}{}
\titleformat{\subsection}{\normalfont\large\bfseries\color{hdr}}{\thesubsection}{0.6em}{}

\pagestyle{fancy}
\fancyhf{}
\fancyhead[L]{\small\textit{Guía de comprensión de datos -- CICIDS-2018}}
\fancyhead[R]{\small\thepage}
\renewcommand{\headrulewidth}{0.4pt}

\newcommand{\code}[1]{\texttt{\small #1}}

\newtcolorbox{cbox}[2]{colback=#1!5,colframe=#1,fonttitle=\bfseries,title={#2},
  boxrule=0.9pt,left=5pt,right=5pt,top=4pt,bottom=4pt}

\title{\textbf{\Large Guía de comprensión de datos}\\[4pt]
\Large \textit{Qué es CICIDS-2018, cómo se lee, y qué te va a morder}\\[8pt]
\normalsize Para alguien que nunca vio este dataset}
\author{Repositorio \code{Ironsss/iag-symmetry}}
\date{2026-08-29}

\begin{document}
\maketitle
\thispagestyle{empty}

\begin{cbox}{hdr}{Para quién es esto}
Este documento asume que no conocés CICIDS-2018. No busca hallazgos de investigación
--- eso está en \code{docs/datos-reales/ANALISIS-DATOS-REALES.pdf} --- busca que puedas
abrir un archivo de \code{data/raw/} y entender qué estás mirando sin tener que reconstruir
el contexto vos mismo. Cada número y figura viene de los dos archivos reales que este
repositorio ya tiene descargados, vía \code{scripts/analyze\_real\_data.py} y
\code{scripts/explore\_real\_data\_structure.py}.
\end{cbox}

\tableofcontents
\newpage

% ============================================================
\section{Qué es CICIDS-2018}
% ============================================================

CICIDS-2018 (Sharafaldin, Lashkari \& Ghorbani, 2018 --- \code{sharafaldin2018cicids} en
\code{paper/references.bib}, verificada) es un dataset público de tráfico de red, generado
en un testbed real (no simulado) con tráfico benigno de usuarios reales y ataques
ejecutados a propósito, en días separados, contra una infraestructura AWS. Cada día del
dataset corresponde a un tipo de ataque (o a tráfico puramente benigno): el archivo que
usa este análisis, \code{Thuesday-20-02-2018}, es el día de DDoS-LOIC-HTTP; el otro,
\code{Thursday-01-03-2018}, es el día de Infiltración.

El tráfico crudo (paquetes) se procesó con \textbf{CICFlowMeter}, una herramienta que
reconstruye cada conversación de red paquete por paquete y la resume en \textbf{una fila
de estadísticas} por conversación. Ese resumen --- no los paquetes --- es lo que se
distribuye públicamente y lo que este repositorio consume. La etiqueta (\code{Label}) no
viene de un modelo ni de una heurística: viene del \textbf{calendario conocido de
ataques} del testbed --- se sabe qué IP atacó a qué IP y cuándo, porque quien generó el
dataset lo orquestó así.

% ============================================================
\section{Qué es un ``flujo''}
% ============================================================

Un \textbf{flujo} (flow) es una conversación entre dos hosts: mismo par de IPs, mismo par
de puertos, mismo protocolo, dentro de una ventana de tiempo continua. \textbf{Una fila
del CSV = un flujo}, no un paquete --- un flujo típico contiene entre unos pocos y miles
de paquetes, todos resumidos (promedio, máximo, mínimo, desviación estándar) en las
columnas de esa fila.

\begin{cbox}{good}{La distinción que más confunde al empezar}
\code{Tot Fwd Pkts} no es ``cuántos paquetes hay en el dataset'' --- es ``cuántos
paquetes tuvo \emph{este flujo en particular}, en la dirección del que inició la
conexión (forward)''. Cada fila es autocontenida: sus propios conteos, sus propios
tiempos, su propia etiqueta.
\end{cbox}

``Forward'' (\code{Fwd}) es la dirección del host que abrió la conexión hacia el otro;
``backward'' (\code{Bwd}) es la respuesta. En una conexión web normal, el cliente
manda poco (una petición HTTP) y el servidor manda mucho (la página) --- \code{Bwd} suele
ser mayor que \code{Fwd}. Guardá esa intuición: en \S\ref{sec:gotchas} vas a ver un
ataque real donde pasa exactamente eso, al revés de lo que uno esperaría de un
``atacante''.

% ============================================================
\section{Las familias de columnas, en lenguaje llano}
% ============================================================

84 columnas crudas (78 si el archivo no trae IPs --- ver \S\ref{sec:gotchas}), pero se
organizan en pocas familias. La lista completa columna-por-columna con su nombre
normalizado está en \code{docs/datos-reales/ANALISIS-DATOS-REALES.pdf} \S2; acá va el
``para qué sirve cada familia'':

\begin{center}
\begin{longtable}{@{}lp{10cm}@{}}
\toprule
\textbf{Familia} & \textbf{Qué responde} \\
\midrule
\endhead
Identificación & ¿Quién habló con quién? \code{Src/Dst IP}, \code{Src/Dst Port},
  \code{Protocol} (6 = TCP, 17 = UDP). \code{Flow ID} es sólo la concatenación de las
  cuatro anteriores, un identificador, no una feature. \\
Temporales & ¿Cuándo, y durante cuánto tiempo? \code{Timestamp} (inicio del flujo),
  \code{Flow Duration}, y las familias \code{IAT} (Inter-Arrival Time: tiempo entre
  paquetes consecutivos, por dirección y total) --- captura si el flujo fue una ráfaga
  rápida o una conversación espaciada. \\
Volumen & ¿Cuánto se mandó? \code{Tot Fwd/Bwd Pkts} (conteo de paquetes),
  \code{TotLen Fwd/Bwd Pkts} (bytes totales), y las tasas derivadas
  \code{Flow Byts/s}, \code{Flow Pkts/s}. \\
Forma del paquete & ¿Los paquetes eran todos iguales o muy variables?
  \code{Fwd/Bwd Pkt Len Max/Min/Mean/Std}, y sus versiones sin separar dirección
  (\code{Pkt Len *}). Una ráfaga de paquetes idénticos (std $\approx$0) se ve muy distinta
  de tráfico web mixto. \\
Banderas TCP & ¿Qué tipo de paquetes TCP hubo? \code{FIN/SYN/RST/PSH/ACK/URG Flag Cnt}
  --- conteos de banderas de control TCP; patrones anómalos (muchos \code{SYN} sin
  \code{ACK}, por ejemplo) son la firma clásica de varios ataques de red, aunque no de
  los que este análisis mide. \\
Ratios derivados & \code{Down/Up Ratio} (bytes bajada/subida), \code{Pkt Size Avg},
  \code{Fwd/Bwd Seg Size Avg} --- resúmenes de las familias anteriores, pensados para
  ahorrarle trabajo a un modelo que no puede ver la distribución completa. \\
Actividad/inactividad & \code{Active Mean/Std/Max/Min}, \code{Idle Mean/Std/Max/Min} ---
  cuánto tiempo el flujo estuvo ``activo'' (paquetes seguidos) vs.\ ``inactivo'' (pausas
  largas) dentro de su propia duración. \\
Ventana inicial & \code{Init Fwd/Bwd Win Byts} --- tamaño de ventana TCP inicial;
  \textbf{usa $-1$ como centinela de ``desconocido''}, no como valor real (ver
  \S\ref{sec:gotchas}). \\
Etiqueta & \code{Label} --- texto libre (\code{Benign}, \code{DDoS attacks-LOIC-HTTP},
  \code{Infilteration}, etc.), la verdad de campo del testbed. \\
\bottomrule
\end{longtable}
\end{center}

% ============================================================
\section{Una fila real, anotada}
% ============================================================

Un flujo de ataque real, tal cual aparece en \code{Thuesday-20-02-2018\_TrafficForML\_CICFlowMeter.csv}
(vía \code{scripts/explore\_real\_data\_structure.py::annotated\_sample\_flow}):

\begin{center}
\begin{tabular}{@{}lll@{}}
\toprule
\textbf{Columna} & \textbf{Valor} & \textbf{Lectura} \\
\midrule
\code{Src IP : Src Port} & 18.216.200.189:65404 & el atacante \\
\code{Dst IP : Dst Port} & 172.31.69.25:80 & la víctima, puerto HTTP \\
\code{Timestamp} & 2018-02-20 10:44:22 & dentro de la ventana de ataque real \\
\code{Flow Duration} & 1.382.540 $\mu$s & $\approx$1.38 segundos \\
\code{Tot Fwd/Bwd Pkts} & 3 / 4 & pocos paquetes en ambas direcciones \\
\code{TotLen Fwd/Bwd Pkts} & \textbf{20 / 964 bytes} & el atacante mandó 20 bytes, recibió 964 \\
\code{Label} & DDoS attacks-LOIC-HTTP & confirmado como ataque \\
\bottomrule
\end{tabular}
\end{center}

\begin{cbox}{warn}{Por qué esta fila importa más de lo que parece}
\code{TotLen Fwd Pkts} (20) $\ll$ \code{TotLen Bwd Pkts} (964): el atacante manda poco y
\emph{recibe} más --- exactamente lo opuesto a la intuición de ``un atacante inunda con
tráfico''. LOIC-HTTP manda peticiones HTTP pequeñas y recibe la respuesta (o página de
error) del servidor. Es la razón concreta, en un flujo real, de por qué
\code{docs/datos-reales/ANALISIS-DATOS-REALES.pdf} mide un AUC de asimetría per-flujo
\textbf{invertido} (0.381, peor que el azar) en este mismo archivo: la señal
``manda-más-de-lo-que-recibe'' no apunta hacia el ataque en este tipo de tráfico, apunta
en contra.
\end{cbox}

% ============================================================
\section{Balance de clases}
% ============================================================

\begin{center}
\begin{tabular}{@{}lrrr@{}}
\toprule
\textbf{Archivo} & \textbf{Filas} & \textbf{\% ataque} & \textbf{Etiquetas presentes} \\
\midrule
Thursday-01-03-2018 & 331.126 & 28.1\% & Benign, Infilteration \\
Thuesday-20-02-2018 & 7.948.746 & 7.2\% & Benign, DDoS attacks-LOIC-HTTP \\
\bottomrule
\end{tabular}
\end{center}

7.2\% de positivos (DDoS) es un desbalanceo real, no extremo --- suficiente para que un
modelo ingenuo (``predecir siempre benigno'') ya acierte 92.8\% de las veces, sin
detectar un solo ataque. Es el régimen para el que técnicas como SMOTE (Chawla et al.\
2002, ya citada en \code{paper/references.bib}) existen; ver el análisis de clustering en
\code{docs/datos-reales/} \S6 para cómo se ve este desbalanceo en la estructura de los
datos, no sólo en el conteo.

% ============================================================
\section{El grafo de interacciones}
% ============================================================

Cada flujo conecta dos hosts. Si se dibuja un nodo por host y una arista por cada par que
conversó, aparece el grafo de interacciones de la red --- el objeto que le da nombre al
proyecto (IAG = índice de asimetría de \emph{grafos}).

\begin{figure}[ht]
\centering
\includegraphics[width=0.8\textwidth]{../../results/figures/real_data_interaction_graph_attack.png}
\caption{Una ventana real de 60 segundos con ataque. Cada punto es un host; cada línea, un
flujo. El grupito rojo abajo a la derecha son los $\sim$10 hosts atacantes hablándole a la
víctima --- notá que no toca al resto del gráfico: el ataque, en esta ventana, es un
\emph{componente aparte}, no está mezclado con el tráfico normal.}
\end{figure}

Leerlo: el tamaño de cada punto es cuántas conexiones tiene ese host (su \emph{grado});
una flecha va de quien inició la conexión a quien la recibió. La estructura ``un
puñado de hosts muy conectados + muchos hosts con pocas conexiones'' (visible en la
ventana benigna, \code{docs/datos-reales/ANALISIS-DATOS-REALES.pdf} fig.\ 5) es típica de
tráfico normal cliente-servidor; un grupo de puntos aislado, hablándole todos a un mismo
destino, es la firma visual de un ataque de denegación de servicio.

% ============================================================
\section{Los gotchas reales de este dataset}
\label{sec:gotchas}
% ============================================================

Encontrados corriendo el código contra los archivos reales en esta sesión --- no son
advertencias genéricas de ``ojo con los datos sucios'', son cosas específicas de
CICIDS-2018 que rompieron código real de este repositorio hasta que se corrigieron:

\begin{center}
\begin{longtable}{@{}p{4.2cm}p{9.5cm}@{}}
\toprule
\textbf{Gotcha} & \textbf{Qué hace} \\
\midrule
\endhead
Esquema no uniforme & La mayoría de los archivos ``Processed Traffic Data'' empiezan en
  \code{Dst Port} --- \textbf{sin} \code{Src IP}/\code{Dst IP}. De los 10 archivos del
  bucket público, sólo uno (\code{Thuesday-20-02-2018}) trae el esquema ancho con IPs. Sin
  IPs no hay grafo posible. \\
Headers duplicados incrustados & Algunos archivos son la concatenación de varios exports
  de CICFlowMeter, cada uno con su propia fila de encabezado repetida en medio del
  archivo. Sin cuidado, esto corrompe la inferencia de tipos de cualquier parser (ver
  \code{docs/datos-reales/ANALISIS-DATOS-REALES.pdf} \S3) --- una columna entera puede
  leerse como texto en vez de número, silenciosamente. \\
\code{Infinity} literal & \code{Flow Byts/s} y \code{Flow Pkts/s} dividen por
  \code{Flow Duration}; cuando la duración es 0, CICFlowMeter escribe la palabra
  \code{Infinity}, no un valor faltante. Hay que decidir explícitamente qué hacer con
  eso (este repositorio lo reemplaza por 0). \\
Centinela $-1$ & \code{Init Fwd/Bwd Win Byts} usa $-1$ para ``no se pudo determinar'',
  no como una medida real negativa. Cualquier transformación que asuma valores no
  negativos (un logaritmo, por ejemplo) tiene que tratar esto aparte. \\
Cola pesada + inflación en cero & La mayoría de las columnas de conteo/flags son 0 para
  la inmensa mayoría de los flujos, pero unos pocos flujos tienen valores enormes
  (duración de horas, millones de bytes). Un escalado ingenuo (media/desvío) deja que
  esos pocos flujos dominen cualquier análisis multivariado --- pasó de verdad al
  intentar proyectar estos datos a 2D (\code{docs/datos-reales/ANALISIS-DATOS-REALES.pdf}
  \S6). \\
\bottomrule
\end{longtable}
\end{center}

\end{document}
